%------------------------------------------------------------------------------------
%  SECTION 1.2 : LE MOTEUR LUALATEX ET LA GESTION MODERNE DES POLICES
%------------------------------------------------------------------------------------
\section{Moteur Lua\LaTeX{} et Gestion Dynamique des Polices}
\label{sec:lualatex_polices}

Pendant plusieurs décennies, le compilateur historique \textsf{pdf\LaTeX{}} a dominé le paysage éditorial.
Cependant, ses limitations intrinsèques liées à l'encodage sur 8 bits nécessitaient une armada de paquets intermédiaires (\texttt{inputenc}, \texttt{fontenc}, \texttt{textcomp}) et rendaient l'utilisation des polices vectorielles modernes OpenType particulièrement laborieuse.

Le modèle \texttt{ClasseGuillaume} s'appuie résolument et exclusivement sur \textbf{Lua\LaTeX{}}, successeur moderne alliant un moteur \TeX{} natif Unicode à la puissance du langage de script embarqué Lua.

\subsection{L'abandon des paquets d'encodage obsolètes}
\label{subsec:unicode_natif}

Sous Lua\LaTeX{}, le code source \LaTeX{} est traité nativement en \textbf{UTF-8}.
Les instructions suivantes, jadis ubiquitaires, sont devenues \textbf{strictement obsolètes} et ne doivent jamais figurer dans votre document :
\begin{itemize}
  \item \texttt{\textbackslash usepackage[utf8]\{inputenc\}} : Inutile et néfaste, Lua\LaTeX{} gère l'Unicode dès sa première ligne.
  \item \texttt{\textbackslash usepackage[T1]\{fontenc\}} : Obsolète, le système de polices \textsf{fontspec} régit directement le rendu glyphe par glyphe.
\end{itemize}

Grâce à ce traitement natif, tous les caractères accentués français (é, è, à, ç, œ, « ») ainsi que les symboles scientifiques et alphabets internationaux sont insérés sans aucune macro d'échappement.

\subsection{Les trois modes de chargement des polices typographiques}
\label{subsec:trois_modes_polices}

Une question cruciale lors de la transition entre une plateforme cloud comme Overleaf et une installation locale (TeX Live) concerne l'accessibilité des polices de caractères :

\begin{enumerate}
  \item \textbf{Mode 1 : Les polices système OS (la méthode la plus directe) :}
  Sur une machine équipée de Windows ou macOS, le sous-système \textsf{luaotfload} indexe automatiquement le dossier des polices système (\texttt{C:\textbackslash Windows\textbackslash Fonts\textbackslash}).
  N'importe quelle police installée sur votre ordinateur (telle que \textsf{Calibri}, \textsf{Segoe UI}, \textsf{Arial}, \textsf{Consolas}) peut être appelée directement par son nom usuel :
  \begin{verbatim}
  \setmainfont{Calibri}
  \setsansfont{Calibri}
  \end{verbatim}
  Le compilateur associe automatiquement les variantes grasses (\texttt{calibrib.ttf}), italiques (\texttt{calibrii.ttf}) et grasses-italiques sans aucune intervention manuelle.
  Aucun fichier de police n'a besoin d'être copié dans le projet !

  \item \textbf{Mode 2 : Les polices OpenType natives de \TeX{} Live (universelles et libres) :}
  La distribution complète \TeX{} Live installe sur votre machine des centaines de familles de polices professionnelles prêtes à l'emploi.
  Par exemple, la suite \textsf{\TeX{} Gyre} fournit des équivalents modernes et fidèles des plus grands classiques :
  \begin{itemize}
    \item \textsf{TeX Gyre Heros} (équivalent moderne d'Helvetica),
    \item \textsf{TeX Gyre Termes} (équivalent de Times),
    \item \textsf{TeX Gyre Pagella} (équivalent de Palatino),
    \item \textsf{Libertinus Serif} et \textsf{Fira Sans}.
  \end{itemize}
  Ces polices fonctionnent de manière identique et immédiate sur Windows, Linux, macOS et Overleaf, sans aucun fichier physique dans le répertoire du projet.

  \item \textbf{Mode 3 : Le repli local dans le projet (\texttt{./Fonts/...}) (méthode Overleaf) :}
  Sur les serveurs cloud comme Overleaf (qui tournent sous des conteneurs Linux), les polices propriétaires Microsoft ne sont pas préinstallées sous licence libre.
  Pour pallier cette absence sans faire échouer la compilation, \texttt{ClasseGuillaume} embarque un dossier local de secours \texttt{./Fonts/}.
  Grâce à la commande conditionnelle intelligente \texttt{\textbackslash IfFontExistsTF}, la classe interroge d'abord le système de votre PC ; si la police y est absente (ex. sur Overleaf), elle active automatiquement le fichier de repli local.
\end{enumerate}

\subsection{Options de classe disponibles pour le choix de la police}
\label{subsec:options_polices}

Dans l'en-tête de votre document principal \texttt{main.tex}, vous pouvez sélectionner votre famille de polices en un mot-clé au sein de la commande \texttt{\textbackslash documentclass} :

\begin{itemize}
  \item \texttt{font=calibri} *(option par défaut)* : Active \textsf{Calibri} pour un style sobre, fluide et moderne, hautement lisible sur écran comme à l'impression.
  \item \texttt{font=heros} : Active \textsf{TeX Gyre Heros}, conférant un rendu suisse net et contemporain.
  \item \texttt{font=libertinus} : Active \textsf{Libertinus Serif}, idéale pour les mémoires littéraires ou les thèses académiques traditionnelles à fort contenu textuel.
  \item \texttt{font=comicsans} : Active \textsf{Comic Sans MS} (utile à des fins pédagogiques ou pour des présentations informelles).
  \item \texttt{font=standard} : Active la police standard \textsf{Latin Modern}.
\end{itemize}

\subsection{Démonstration des commandes de polices locales}
\label{subsec:commandes_polices_locales}

Outre la police générale du document, la classe met à votre disposition des commandes spécifiques pour des encarts, titres ou spécimens ponctuels :

\begin{itemize}
  \item \textbf{Police de titre de chapitre :} La commande \texttt{\textbackslash AsapFont} applique la police {\AsapFont Asap Condensed}, conférant un impact visuel fort aux intitulés de chapitres.
  \item \textbf{Spécimen Calibri :} Accessible ponctuellement via {\CalibriFont \textbackslash CalibriFont pour forcer le style sans-serif}.
  \item \textbf{Spécimen fantaisie :} Accessible via {\ComicSansMSFont \textbackslash ComicSansMSFont pour des démonstrations ciblées}.
\end{itemize}

